\section{Результати}
\label{sec:results}

У цьому розділі порівнюємо \Full{} і \TS{}
за часовими характеристиками поширення інфекції
та допустимістю міжспільнотних переходів.
Для цього проведемо серії експериментів, описані у Розділі~\ref{sec:design}.
Спочатку розглядаємо перший експорт інфекції,
щоб визначити наслідки втрати інформації
про вузли-мости. Далі досліджуємо поширення
вздовж ланцюга спільнот, вплив характеристик
мережі та зв'язок розбіжностей зі співвідношенням
часових масштабів.

Порівняння ланцюгів охоплює 48 конфігурацій
із 400 парами запусків \Full{} і \TS{} у кожній,
загалом 19\,200 пар.


\subsection{Часова точність і допустимість
міжспільнотного передавання}
\label{sec:res_clique}

Перший експеримент перевіряє, чи може \TS{}
правильно відтворювати час експорту інфекції,
не зберігаючи інформації про конкретні
вузли-мости. Для цього розглядаємо мережу
«кліка--листок», у якій внутрішнє поширення
відбувається швидко.

У цьому випадку розподіли часу експорту
були близькими (рисунок~\ref{fig:clique_leaf}).
Середній час зараження листка становив
110.18 у \Full{} та 112.62 у \TS{},
а нормована відстань Вассерштайна
дорівнювала $\eW=2.95\%$.
Парна різниця середніх становила $+2.43$
із 95\% бутстреп-інтервалом $[-0.35,5.27]$,
який містить нуль. Отже, у цьому експерименті
не виявлено значного систематичного
зміщення середнього часу експорту.

Близькість часових розподілів у цьому
експерименті узгоджується з швидким
і рівномірним поширенням інфекції
всередині кліки. Завдяки її симетричній
структурі зараження швидко охоплює
вузли незалежно від їхнього розташування,
зокрема й вузли-мости.

Очікуваний час зараження першого моста
становить лише $\mathbb{E}[T_B]=8.36$,
тоді як середній час експорту перевищує 110.
Отже, відмінності між фактичною
та агрегованою інтенсивностями передавання
зберігаються протягом порівняно короткого
часу відносно загального очікування експорту,
що узгоджується з невеликою спостереженою
розбіжністю між моделями.

Проте близькість часових розподілів
не означає збереження допустимого
маршруту зараження.
У 98 із 10\,000 запусків \TS{} (0.98\%)
листок було заражено раніше за будь-який
із двох справжніх вузлів-мостів.
У \Full{} такий порядок подій неможливий.
Таким чином, агрегований інтерфейс може
правильно відтворювати часові характеристики
передавання, одночасно допускаючи
переходи, неможливі у початковій мережі.

\begin{figure}[tbp]
\centering
% BEGIN EDITABLE FIGURE clique_leaf
% Editable PGFPlots version of clique_leaf_endpoint_and_mechanism.
\begin{tikzpicture}
\begin{groupplot}[
  group style={group size=3 by 1, horizontal sep=1.55cm},
  paper axis, width=.24\linewidth, height=3.7cm, scale only axis,
  legend style={font=\scriptsize, draw=none},
]
\nextgroupplot[
  title={\textbf{a} Export times},
  xlabel={total export time}, ylabel={empirical CDF},
  xmin=0, xmax=600, ymin=0, ymax=1.01,
  legend style={at={(.98,.36)}, anchor=south east},
]
  \addplot+[color=FullBlue, solid, no markers, const plot, line width=.8pt]
    table[x=x,y=y]{figures_pgfplots/leaf_clique_export_full.dat};
  \addlegendentry{Full}
  \addplot+[color=ScaleOrange, dashed, no markers, const plot, line width=.8pt]
    table[x=x,y=y]{figures_pgfplots/leaf_clique_export_twoscale.dat};
  \addlegendentry{TwoScale}
  \node[anchor=south east, align=right, fill=white, draw=black!20, inner sep=2pt, font=\scriptsize]
    at (rel axis cs:.97,.07) {$W_1=3.249$\\$D_{\rm KS}=0.0142$};
\nextgroupplot[
  title={\textbf{b} Export minus bridge},
  xlabel={$T_{\mathrm{export}}-T_B$}, ylabel={empirical CDF},
  xmin=-20, xmax=100, ymin=0, ymax=.55,
]
  \path[fill=black!8] (axis cs:-20,0) rectangle (axis cs:0,.55);
  \addplot+[black!75, no markers, line width=.35pt] coordinates {(0,0) (0,.55)};
  \addplot+[color=FullBlue, solid, no markers, const plot, line width=.8pt]
    table[x=x,y=y]{figures_pgfplots/leaf_clique_signed_delay_full.dat};
  \addplot+[color=ScaleOrange, dashed, no markers, const plot, line width=.8pt]
    table[x=x,y=y]{figures_pgfplots/leaf_clique_signed_delay_twoscale.dat};
  \node[anchor=north west, font=\scriptsize, text=black!65] at (axis cs:-18,.51) {export before bridge};
\nextgroupplot[
  title={\textbf{c} Bridge bypass},
  ylabel={bridge-bypass runs (\%)},
  xmin=.5, xmax=2.5, ymin=-.03, ymax=1.35, clip=false,
  xtick={1,2}, xticklabels={Full,TwoScale},
]
  \addplot+[
    only marks, mark=*, mark size=1.7pt, color=black!80,
    error bars/.cd, y dir=both, y explicit,
  ] table[x=x,y=rate_pct,
      y error minus expr=\thisrow{rate_pct}-\thisrow{low_pct},
      y error plus expr=\thisrow{high_pct}-\thisrow{rate_pct}]
      {figures_pgfplots/leaf_clique_bypass.dat};
  \node[anchor=south, align=center, font=\scriptsize] at (axis cs:1,.14) {0/10000\\(0.00\%)};
  \node[anchor=south, align=center, font=\scriptsize] at (axis cs:2,1.18) {98/10000\\(0.98\%)};
\end{groupplot}
\end{tikzpicture}
% END EDITABLE FIGURE clique_leaf
\caption{%
Перший експорт інфекції в мережі
«кліка--листок» ($n=100$, $k=2$,
$\beta=0.005$, 10\,000 запусків кожної моделі).
(a) Емпіричні функції розподілу часу експорту.
(b) Розподіли різниці $T_{\mathrm{export}}-T_B$;
від'ємні значення відповідають експорту
до зараження будь-якого справжнього вузла-моста.
(c) Частка запусків із таким порушенням
порядку передавання та 95\% довірчі інтервали.
}
\label{fig:clique_leaf}
\end{figure}

Щоб з'ясувати, як цей ефект залежить
від внутрішньої структури, замінюємо
кліку випадковими зв'язними графами
зі змінною середньою довжиною
найкоротших шляхів (ASPL).
У цій серії зі збільшенням ASPL
також зменшується щільність графів.

У щільних мережах із короткими шляхами
передчасні експорти залишалися рідкісними:
їхня частка не перевищувала 0.5\%
за ASPL $\leq 1.75$.
Зі збільшенням ASPL вона зростала
і за ASPL 6 досягала 32.5\%
для обох досліджених значень $k$
(рисунок~\ref{fig:source_sweep}).

Водночас змінювалися й часові характеристики
експорту. За ASPL 6 середній час експорту
у \TS{} був меншим, ніж у \Full{},
на 8.3\% для $k=1$ і на 6.2\% для $k=5$.
Відповідні 95\% бутстреп-інтервали
становили $[-13.1,-3.5]\%$
та $[-10.3,-2.3]\%$.
У жодній конфігурації цієї серії
не спостерігалося додатного
зсуву середнього часу експорту.

Цей напрям зміщення узгоджується
з твердженням~\ref{prop:ordering}
для ідеального двомасштабного процесу:
усереднення інтенсивностей за випадковими
вузлами-мостами може спричиняти
систематично раніший експорт.

\begin{figure}[tbp]
\centering
% BEGIN EDITABLE FIGURE source_sweep
% Editable PGFPlots version of fig_aspl_export_fidelity.
\begin{tikzpicture}
\begin{groupplot}[
  group style={group size=3 by 2, horizontal sep=1.55cm, vertical sep=1.35cm},
  paper axis, width=.24\linewidth, height=3.7cm, scale only axis,
  xmin=.85, xmax=6.15, xtick={1,2,...,6},
]
\nextgroupplot[
  title={$k=1$: export time}, ylabel={mean export time},
  xticklabels=\empty, ymin=0, ymax=1643.82646974,
  legend style={font=\scriptsize, draw=none, at={(.98,.98)}, anchor=north east},
]
  \addplot+[color=FullBlue, mark=*, mark size=1.1pt, mark options={fill=FullBlue,draw=FullBlue}, line width=.6pt,
    error bars/.cd, y dir=both, y explicit]
    table[x=aspl,y=full_mean,y error minus expr=\thisrow{full_mean}-\thisrow{full_low},y error plus expr=\thisrow{full_high}-\thisrow{full_mean}]
    {figures_pgfplots/leaf_source_sweep_k1.dat};
  \addlegendentry{Full}
  \addplot+[color=ScaleOrange, mark=square*, mark size=1.1pt, mark options={fill=ScaleOrange,draw=ScaleOrange}, line width=.6pt,
    error bars/.cd, y dir=both, y explicit]
    table[x=aspl,y=twoscale_mean,y error minus expr=\thisrow{twoscale_mean}-\thisrow{twoscale_low},y error plus expr=\thisrow{twoscale_high}-\thisrow{twoscale_mean}]
    {figures_pgfplots/leaf_source_sweep_k1.dat};
  \addlegendentry{TwoScale}
\nextgroupplot[title={$k=1$: $W_1$ error}, ylabel={$W_1/\overline T_{\rm Full}$ (\%)}, xticklabels=\empty, ymin=0, ymax=13.026367574]
  \addplot+[color=ScaleOrange, mark=*, mark size=1.1pt, mark options={fill=ScaleOrange,draw=ScaleOrange}, line width=.6pt]
    table[x=aspl,y=w1]{figures_pgfplots/leaf_source_sweep_k1.dat};
\nextgroupplot[title={$k=1$: bridge bypass}, ylabel={export before any bridge (\%)}, xticklabels=\empty, ymin=0, ymax=39.9]
  \addplot+[color=ScaleOrange, mark=*, mark size=1.1pt, mark options={fill=ScaleOrange,draw=ScaleOrange}, line width=.6pt]
    table[x=aspl,y=bypass]{figures_pgfplots/leaf_source_sweep_k1.dat};
  \addplot+[color=ScaleOrange, mark=none, line width=.45pt,
    error bars/.cd, y dir=both, y explicit]
    table[x=aspl,y=bypass,y error minus expr=\thisrow{bypass}-\thisrow{bypass_low},y error plus expr=\thisrow{bypass_high}-\thisrow{bypass}]
    {figures_pgfplots/leaf_source_sweep_k1_bypass_ci.dat};
\nextgroupplot[title={$k=5$: export time}, xlabel={target ASPL}, ylabel={mean export time}, ymin=0, ymax=1643.82646974]
  \addplot+[color=FullBlue, mark=*, mark size=1.1pt, mark options={fill=FullBlue,draw=FullBlue}, line width=.6pt,
    error bars/.cd, y dir=both, y explicit]
    table[x=aspl,y=full_mean,y error minus expr=\thisrow{full_mean}-\thisrow{full_low},y error plus expr=\thisrow{full_high}-\thisrow{full_mean}]
    {figures_pgfplots/leaf_source_sweep_k5.dat};
  \addplot+[color=ScaleOrange, mark=square*, mark size=1.1pt, mark options={fill=ScaleOrange,draw=ScaleOrange}, line width=.6pt,
    error bars/.cd, y dir=both, y explicit]
    table[x=aspl,y=twoscale_mean,y error minus expr=\thisrow{twoscale_mean}-\thisrow{twoscale_low},y error plus expr=\thisrow{twoscale_high}-\thisrow{twoscale_mean}]
    {figures_pgfplots/leaf_source_sweep_k5.dat};
\nextgroupplot[title={$k=5$: $W_1$ error}, xlabel={target ASPL}, ylabel={$W_1/\overline T_{\rm Full}$ (\%)}, ymin=0, ymax=13.026367574]
  \addplot+[color=ScaleOrange, mark=*, mark size=1.1pt, mark options={fill=ScaleOrange,draw=ScaleOrange}, line width=.6pt]
    table[x=aspl,y=w1]{figures_pgfplots/leaf_source_sweep_k5.dat};
\nextgroupplot[title={$k=5$: bridge bypass}, xlabel={target ASPL}, ylabel={export before any bridge (\%)}, ymin=0, ymax=39.9]
  \addplot+[color=ScaleOrange, mark=*, mark size=1.1pt, mark options={fill=ScaleOrange,draw=ScaleOrange}, line width=.6pt]
    table[x=aspl,y=bypass]{figures_pgfplots/leaf_source_sweep_k5.dat};
  \addplot+[color=ScaleOrange, mark=none, line width=.45pt,
    error bars/.cd, y dir=both, y explicit]
    table[x=aspl,y=bypass,y error minus expr=\thisrow{bypass}-\thisrow{bypass_low},y error plus expr=\thisrow{bypass_high}-\thisrow{bypass}]
    {figures_pgfplots/leaf_source_sweep_k5_bypass_ci.dat};
\end{groupplot}
\end{tikzpicture}
% END EDITABLE FIGURE source_sweep
\caption{%
Вплив внутрішньої структури джерела
на перший експорт до листка
($n=1000$, $\beta=0.003$, $k=1$ угорі
та $k=5$ унизу; 400 парних запусків
для кожної конфігурації).
Ліворуч: середній час експорту.
У центрі: нормована відстань Вассерштайна
$\eW$ між розподілами часу експорту.
Праворуч: частка експортів \TS{}
до зараження першого справжнього моста;
у \Full{} такі події неможливі.
Відрізки похибок --- 95\% інтервали
блокового бутстрепу.
}
\label{fig:source_sweep}
\end{figure}

Оцінюючи відстань Вассерштайна,
потрібно враховувати вибіркову мінливість.
Для майже експоненційних часів експорту
за $R=400$ очікувана нормована відстань
між двома незалежними вибірками
з одного розподілу становить приблизно
$\sqrt{\pi/R}\approx9\%$
(Додаток~\ref*{app:floor}).
Тому невеликі значення $\eW$
в одиничному експорті самі по собі
ще не свідчать про відмінність моделей.

На відміну від часової відстані,
кожен передчасний експорт безпосередньо
вказує на порушення допустимого порядку
передавання. Наприклад, за ASPL 4
і $k=1$ такі порушення спостерігалися
у 19.75\% запусків, хоча часові
розподіли залишалися порівняно близькими.

Отже, перший експорт демонструє
два різні наслідки агрегації.
Навіть за близьких часів передавання
модель може порушувати фізичні
обмеження міжспільнотних контактів.
Коли внутрішнє поширення стає повільнішим
і нерівномірнішим, ця втрата інформації
може також спричиняти помітне
зміщення часу експорту.


\subsection{Поширення розбіжностей
уздовж ланцюга спільнот}
\label{sec:res_chain_random}

Одиничний експорт не показує,
як відмінності між моделями поводяться
після кількох міжспільнотних переходів.
Для цього досліджуємо поширення інфекції
вздовж ланцюга з десяти спільнот.

Спочатку розглядаємо ланцюг клік,
у якому внутрішнє зараження відбувається
значно швидше за міжспільнотне передавання.
Навіть після дев'яти послідовних переходів
часові розподіли залишалися близькими
(рисунок~\ref{fig:chain_cliques}).

Медіанний час прибуття інфекції
до останньої спільноти становив
2785 у \Full{} та 2877 у \TS{}.
Нормована відстань Вассерштайна
дорівнювала 3.34\% і не перевищувала
перестановочного порога 5.1\%
($p=0.32$). Отже, за цього обсягу
вибірки тест не виявив статистично
значущої відмінності між розподілами.

Для кліки з $n$ вузлів за
$\beta_{\mathrm{in}}=\beta_{\mathrm{out}}$
відношення середнього часу зараження
всієї кліки до середнього часу
передавання через одне активне
міжспільнотне ребро становить
$2H_{n-1}/n$, тобто лише близько
$1.5\%$ за $n=1000$.

Отже, внутрішнє поширення значно
швидше за міжспільнотне передавання,
що пояснює невеликі спостережені
розбіжності між моделями. Крім того, інформація про те, які саме вузли є мостами, має порівняно невеликий вплив на час переходу.

\begin{figure}[tbp]
\centering
% BEGIN EDITABLE FIGURE chain_cliques
\begin{tikzpicture}
\begin{groupplot}[paper axis,group style={group size=2 by 1,horizontal sep=1.8cm},width=.38\linewidth,height=4.6cm,scale only axis,xmin=0,xmax=6000,xtick={0,2000,4000,6000},xlabel={Time}]
\nextgroupplot[title={(a) Epidemic propagation},ylabel={Infected nodes},ymin=0,ymax=10300,ytick={0,2500,5000,7500,10000},scaled y ticks=false,legend pos=north west]
\addplot[draw=none,name path=fulllo,forget plot] table[x=time,y=lower]{figures_pgfplots/chain_full_trajectory.dat};
\addplot[draw=none,name path=fullhi,forget plot] table[x=time,y=upper]{figures_pgfplots/chain_full_trajectory.dat};
\addplot[fill=FullBlue,fill opacity=.16,draw=none,forget plot] fill between[of=fulllo and fullhi];
\addplot[FullBlue,thick] table[x=time,y=median]{figures_pgfplots/chain_full_trajectory.dat};
\addlegendentry{\Full{}}
\addplot[draw=none,name path=mmlo,forget plot] table[x=time,y=lower]{figures_pgfplots/chain_mm_trajectory.dat};
\addplot[draw=none,name path=mmhi,forget plot] table[x=time,y=upper]{figures_pgfplots/chain_mm_trajectory.dat};
\addplot[fill=ScaleOrange,fill opacity=.16,draw=none,forget plot] fill between[of=mmlo and mmhi];
\addplot[ScaleOrange,thick] table[x=time,y=median]{figures_pgfplots/chain_mm_trajectory.dat};
\addlegendentry{\TS{}}
\nextgroupplot[title={(b) Arrival in the last community},ylabel={Empirical CDF},ymin=0,ymax=1.02,ytick={0,.25,.5,.75,1}]
\addplot[FullBlue,thick,const plot] table[x=time,y=cdf]{figures_pgfplots/chain_full_cdf.dat};
\addplot[ScaleOrange,thick,const plot] table[x=time,y=cdf]{figures_pgfplots/chain_mm_cdf.dat};
\end{groupplot}
\end{tikzpicture}
% END EDITABLE FIGURE chain_cliques
\caption{%
Поширення інфекції вздовж ланцюга
з десяти клік по 1000 вузлів
($k=1$, $\beta=0.003$, 400 запусків
для кожної моделі).
Ліворуч: медіана кількості інфікованих
вузлів у часі та міжквартильні смуги,
що характеризують розкид між запусками,
а не довірчі інтервали.
Праворуч: емпіричні функції розподілу
часу прибуття до останньої спільноти
$T_9$. Чотири значення \Full{}
і два значення \TS{} лежать поза
показаним часовим діапазоном;
усі вони враховані в аналізі.
}
\label{fig:chain_cliques}
\end{figure}

У ланцюгах випадкових спільнот
із довшими внутрішніми шляхами
спостерігається інша картина.
\TS{} систематично прогнозує
раніше прибуття інфекції,
проте величина цього випередження
змінюється по-різному в абсолютному
та відносному вираженні
(рисунок~\ref{fig:chain_heatmap}).

За ASPL 6 і $k=1$ медіанний
час першого переходу у \TS{}
був меншим на 7.1\%,
а час досягнення останньої спільноти
--- на 13.8\%.
Для $k=5$ відповідні значення
становили 8.7\% та 10.1\%.

Однак відносна розбіжність
не зростала пропорційно кількості
міжспільнотних переходів.
Після початкових змін вона
переважно стабілізувалася
або змінювалася повільно.
Наприклад, за ASPL 5 і $k=1$
медіанне випередження становило
9.6\% для першого переходу
та 8.7\% для останнього.

Водночас абсолютна різниця
між медіанними часами прибуття
в цій конфігурації зросла
приблизно зі 107 до 896
одиниць часу.

Отже, послідовні міжспільнотні
переходи можуть накопичувати
значну абсолютну часову розбіжність,
хоча її відносна величина
залишається приблизно сталою.
Це важливо для задач прогнозування,
у яких навіть відносно невелика
похибка може відповідати значній
різниці в моменті прибуття інфекції
до віддаленої спільноти.

\begin{figure}[tbp]
\centering
% BEGIN EDITABLE FIGURE chain_heatmap
\begin{tikzpicture}
\begin{groupplot}[paper axis,group style={group size=2 by 2,horizontal sep=3.4cm,vertical sep=1.7cm},width=.31\linewidth,height=4.25cm,scale only axis]
\nextgroupplot[title={(a) $k=1$: median shift},xmin=-.5,xmax=8.5,ymin=-.5,ymax=5.5,xtick={0,1,2,3,4,5,6,7,8},xticklabels={1,2,3,4,5,6,7,8,9},ytick={0,1,2,3,4,5},yticklabels={1,2,3,4,5,6},xlabel={Community reached $c$},ylabel={Target internal ASPL},colormap name=paperdiv,point meta min=-14.28295105051232,point meta max=14.28295105051232,colorbar,colorbar style={width=3mm,ylabel={$b_c$ (\%)},ylabel style={font=\scriptsize},scaled y ticks=false},grid=none,tick align=outside]
\addplot[matrix plot*,mesh/cols=9,point meta=explicit] table[x=x,y=y,meta=value]{figures_pgfplots/chain_heatmap_k1_median.dat};
\pgfplotstableread{figures_pgfplots/chain_heatmap_k1_median.dat}\heatmapdata
\pgfplotstableforeachcolumnelement{value}\of\heatmapdata\as\cellvalue{
\pgfplotstablegetelem{\pgfplotstablerow}{x}\of\heatmapdata\edef\cellx{\pgfplotsretval}
\pgfplotstablegetelem{\pgfplotstablerow}{y}\of\heatmapdata\edef\celly{\pgfplotsretval}
\pgfplotstablegetelem{\pgfplotstablerow}{hatch}\of\heatmapdata
\ifnum\pgfplotsretval=1\relax\edef\drawcellhatch{\noexpand\path[pattern=north east lines,pattern color=gray!45] (axis cs:{\cellx-.5},{\celly-.5}) rectangle (axis cs:{\cellx+.5},{\celly+.5});}\drawcellhatch\fi
\pgfmathparse{abs(\cellvalue)>8.2841116092971454}
\ifdim\pgfmathresult pt>0pt\def\celltextcolor{white}\else\def\celltextcolor{black}\fi
\edef\drawcelllabel{\noexpand\node[font=\noexpand\ArrivalCellFont,text=\celltextcolor] at (axis cs:\cellx,\celly) {\noexpand\pgfmathprintnumber[assume math mode=true,fixed,precision=1,fixed zerofill,showpos=false]{\cellvalue}};}\drawcelllabel
}
\draw[white,line width=.4pt] (axis cs:0.5,-.5) -- (axis cs:0.5,5.5);
\draw[white,line width=.4pt] (axis cs:1.5,-.5) -- (axis cs:1.5,5.5);
\draw[white,line width=.4pt] (axis cs:2.5,-.5) -- (axis cs:2.5,5.5);
\draw[white,line width=.4pt] (axis cs:3.5,-.5) -- (axis cs:3.5,5.5);
\draw[white,line width=.4pt] (axis cs:4.5,-.5) -- (axis cs:4.5,5.5);
\draw[white,line width=.4pt] (axis cs:5.5,-.5) -- (axis cs:5.5,5.5);
\draw[white,line width=.4pt] (axis cs:6.5,-.5) -- (axis cs:6.5,5.5);
\draw[white,line width=.4pt] (axis cs:7.5,-.5) -- (axis cs:7.5,5.5);
\draw[white,line width=.4pt] (axis cs:-.5,0.5) -- (axis cs:8.5,0.5);
\draw[white,line width=.4pt] (axis cs:-.5,1.5) -- (axis cs:8.5,1.5);
\draw[white,line width=.4pt] (axis cs:-.5,2.5) -- (axis cs:8.5,2.5);
\draw[white,line width=.4pt] (axis cs:-.5,3.5) -- (axis cs:8.5,3.5);
\draw[white,line width=.4pt] (axis cs:-.5,4.5) -- (axis cs:8.5,4.5);
\nextgroupplot[title={(b) $k=1$: Wasserstein distance},xmin=-.5,xmax=8.5,ymin=-.5,ymax=5.5,xtick={0,1,2,3,4,5,6,7,8},xticklabels={1,2,3,4,5,6,7,8,9},ytick={0,1,2,3,4,5},yticklabels={1,2,3,4,5,6},xlabel={Community reached $c$},ylabel={Target internal ASPL},colormap name=paperseq,point meta min=0,point meta max=13.83377890219645,colorbar,colorbar style={width=3mm,ylabel={$e_W$ (\%)},ylabel style={font=\scriptsize},scaled y ticks=false},grid=none,tick align=outside]
\addplot[matrix plot*,mesh/cols=9,point meta=explicit] table[x=x,y=y,meta=value]{figures_pgfplots/chain_heatmap_k1_w1.dat};
\pgfplotstableread{figures_pgfplots/chain_heatmap_k1_w1.dat}\heatmapdata
\pgfplotstableforeachcolumnelement{value}\of\heatmapdata\as\cellvalue{
\pgfplotstablegetelem{\pgfplotstablerow}{x}\of\heatmapdata\edef\cellx{\pgfplotsretval}
\pgfplotstablegetelem{\pgfplotstablerow}{y}\of\heatmapdata\edef\celly{\pgfplotsretval}
\pgfplotstablegetelem{\pgfplotstablerow}{hatch}\of\heatmapdata
\ifnum\pgfplotsretval=1\relax\edef\drawcellhatch{\noexpand\path[pattern=north east lines,pattern color=gray!45] (axis cs:{\cellx-.5},{\celly-.5}) rectangle (axis cs:{\cellx+.5},{\celly+.5});}\drawcellhatch\fi
\pgfmathparse{abs(\cellvalue)>8.0235917632739397}
\ifdim\pgfmathresult pt>0pt\def\celltextcolor{white}\else\def\celltextcolor{black}\fi
\edef\drawcelllabel{\noexpand\node[font=\noexpand\ArrivalCellFont,text=\celltextcolor] at (axis cs:\cellx,\celly) {\noexpand\pgfmathprintnumber[assume math mode=true,fixed,precision=1,fixed zerofill,showpos=false]{\cellvalue}};}\drawcelllabel
}
\draw[white,line width=.4pt] (axis cs:0.5,-.5) -- (axis cs:0.5,5.5);
\draw[white,line width=.4pt] (axis cs:1.5,-.5) -- (axis cs:1.5,5.5);
\draw[white,line width=.4pt] (axis cs:2.5,-.5) -- (axis cs:2.5,5.5);
\draw[white,line width=.4pt] (axis cs:3.5,-.5) -- (axis cs:3.5,5.5);
\draw[white,line width=.4pt] (axis cs:4.5,-.5) -- (axis cs:4.5,5.5);
\draw[white,line width=.4pt] (axis cs:5.5,-.5) -- (axis cs:5.5,5.5);
\draw[white,line width=.4pt] (axis cs:6.5,-.5) -- (axis cs:6.5,5.5);
\draw[white,line width=.4pt] (axis cs:7.5,-.5) -- (axis cs:7.5,5.5);
\draw[white,line width=.4pt] (axis cs:-.5,0.5) -- (axis cs:8.5,0.5);
\draw[white,line width=.4pt] (axis cs:-.5,1.5) -- (axis cs:8.5,1.5);
\draw[white,line width=.4pt] (axis cs:-.5,2.5) -- (axis cs:8.5,2.5);
\draw[white,line width=.4pt] (axis cs:-.5,3.5) -- (axis cs:8.5,3.5);
\draw[white,line width=.4pt] (axis cs:-.5,4.5) -- (axis cs:8.5,4.5);
\nextgroupplot[title={(c) $k=5$: median shift},xmin=-.5,xmax=8.5,ymin=-.5,ymax=5.5,xtick={0,1,2,3,4,5,6,7,8},xticklabels={1,2,3,4,5,6,7,8,9},ytick={0,1,2,3,4,5},yticklabels={1,2,3,4,5,6},xlabel={Community reached $c$},ylabel={Target internal ASPL},colormap name=paperdiv,point meta min=-14.28295105051232,point meta max=14.28295105051232,colorbar=false,colorbar style={width=3mm,ylabel={$b_c$ (\%)},ylabel style={font=\scriptsize},scaled y ticks=false},grid=none,tick align=outside]
\addplot[matrix plot*,mesh/cols=9,point meta=explicit] table[x=x,y=y,meta=value]{figures_pgfplots/chain_heatmap_k5_median.dat};
\pgfplotstableread{figures_pgfplots/chain_heatmap_k5_median.dat}\heatmapdata
\pgfplotstableforeachcolumnelement{value}\of\heatmapdata\as\cellvalue{
\pgfplotstablegetelem{\pgfplotstablerow}{x}\of\heatmapdata\edef\cellx{\pgfplotsretval}
\pgfplotstablegetelem{\pgfplotstablerow}{y}\of\heatmapdata\edef\celly{\pgfplotsretval}
\pgfplotstablegetelem{\pgfplotstablerow}{hatch}\of\heatmapdata
\ifnum\pgfplotsretval=1\relax\edef\drawcellhatch{\noexpand\path[pattern=north east lines,pattern color=gray!45] (axis cs:{\cellx-.5},{\celly-.5}) rectangle (axis cs:{\cellx+.5},{\celly+.5});}\drawcellhatch\fi
\pgfmathparse{abs(\cellvalue)>8.2841116092971454}
\ifdim\pgfmathresult pt>0pt\def\celltextcolor{white}\else\def\celltextcolor{black}\fi
\edef\drawcelllabel{\noexpand\node[font=\noexpand\ArrivalCellFont,text=\celltextcolor] at (axis cs:\cellx,\celly) {\noexpand\pgfmathprintnumber[assume math mode=true,fixed,precision=1,fixed zerofill,showpos=false]{\cellvalue}};}\drawcelllabel
}
\draw[white,line width=.4pt] (axis cs:0.5,-.5) -- (axis cs:0.5,5.5);
\draw[white,line width=.4pt] (axis cs:1.5,-.5) -- (axis cs:1.5,5.5);
\draw[white,line width=.4pt] (axis cs:2.5,-.5) -- (axis cs:2.5,5.5);
\draw[white,line width=.4pt] (axis cs:3.5,-.5) -- (axis cs:3.5,5.5);
\draw[white,line width=.4pt] (axis cs:4.5,-.5) -- (axis cs:4.5,5.5);
\draw[white,line width=.4pt] (axis cs:5.5,-.5) -- (axis cs:5.5,5.5);
\draw[white,line width=.4pt] (axis cs:6.5,-.5) -- (axis cs:6.5,5.5);
\draw[white,line width=.4pt] (axis cs:7.5,-.5) -- (axis cs:7.5,5.5);
\draw[white,line width=.4pt] (axis cs:-.5,0.5) -- (axis cs:8.5,0.5);
\draw[white,line width=.4pt] (axis cs:-.5,1.5) -- (axis cs:8.5,1.5);
\draw[white,line width=.4pt] (axis cs:-.5,2.5) -- (axis cs:8.5,2.5);
\draw[white,line width=.4pt] (axis cs:-.5,3.5) -- (axis cs:8.5,3.5);
\draw[white,line width=.4pt] (axis cs:-.5,4.5) -- (axis cs:8.5,4.5);
\nextgroupplot[title={(d) $k=5$: Wasserstein distance},xmin=-.5,xmax=8.5,ymin=-.5,ymax=5.5,xtick={0,1,2,3,4,5,6,7,8},xticklabels={1,2,3,4,5,6,7,8,9},ytick={0,1,2,3,4,5},yticklabels={1,2,3,4,5,6},xlabel={Community reached $c$},ylabel={Target internal ASPL},colormap name=paperseq,point meta min=0,point meta max=13.83377890219645,colorbar=false,colorbar style={width=3mm,ylabel={$e_W$ (\%)},ylabel style={font=\scriptsize},scaled y ticks=false},grid=none,tick align=outside]
\addplot[matrix plot*,mesh/cols=9,point meta=explicit] table[x=x,y=y,meta=value]{figures_pgfplots/chain_heatmap_k5_w1.dat};
\pgfplotstableread{figures_pgfplots/chain_heatmap_k5_w1.dat}\heatmapdata
\pgfplotstableforeachcolumnelement{value}\of\heatmapdata\as\cellvalue{
\pgfplotstablegetelem{\pgfplotstablerow}{x}\of\heatmapdata\edef\cellx{\pgfplotsretval}
\pgfplotstablegetelem{\pgfplotstablerow}{y}\of\heatmapdata\edef\celly{\pgfplotsretval}
\pgfplotstablegetelem{\pgfplotstablerow}{hatch}\of\heatmapdata
\ifnum\pgfplotsretval=1\relax\edef\drawcellhatch{\noexpand\path[pattern=north east lines,pattern color=gray!45] (axis cs:{\cellx-.5},{\celly-.5}) rectangle (axis cs:{\cellx+.5},{\celly+.5});}\drawcellhatch\fi
\pgfmathparse{abs(\cellvalue)>8.0235917632739397}
\ifdim\pgfmathresult pt>0pt\def\celltextcolor{white}\else\def\celltextcolor{black}\fi
\edef\drawcelllabel{\noexpand\node[font=\noexpand\ArrivalCellFont,text=\celltextcolor] at (axis cs:\cellx,\celly) {\noexpand\pgfmathprintnumber[assume math mode=true,fixed,precision=1,fixed zerofill,showpos=false]{\cellvalue}};}\drawcelllabel
}
\draw[white,line width=.4pt] (axis cs:0.5,-.5) -- (axis cs:0.5,5.5);
\draw[white,line width=.4pt] (axis cs:1.5,-.5) -- (axis cs:1.5,5.5);
\draw[white,line width=.4pt] (axis cs:2.5,-.5) -- (axis cs:2.5,5.5);
\draw[white,line width=.4pt] (axis cs:3.5,-.5) -- (axis cs:3.5,5.5);
\draw[white,line width=.4pt] (axis cs:4.5,-.5) -- (axis cs:4.5,5.5);
\draw[white,line width=.4pt] (axis cs:5.5,-.5) -- (axis cs:5.5,5.5);
\draw[white,line width=.4pt] (axis cs:6.5,-.5) -- (axis cs:6.5,5.5);
\draw[white,line width=.4pt] (axis cs:7.5,-.5) -- (axis cs:7.5,5.5);
\draw[white,line width=.4pt] (axis cs:-.5,0.5) -- (axis cs:8.5,0.5);
\draw[white,line width=.4pt] (axis cs:-.5,1.5) -- (axis cs:8.5,1.5);
\draw[white,line width=.4pt] (axis cs:-.5,2.5) -- (axis cs:8.5,2.5);
\draw[white,line width=.4pt] (axis cs:-.5,3.5) -- (axis cs:8.5,3.5);
\draw[white,line width=.4pt] (axis cs:-.5,4.5) -- (axis cs:8.5,4.5);
\end{groupplot}
\end{tikzpicture}
% END EDITABLE FIGURE chain_heatmap
\caption{%
Зміна розбіжностей уздовж ланцюга
з десяти спільнот по $n=1000$ вузлів.
Рядки відповідають цільовій ASPL,
стовпці --- досягнутій спільноті
($c=1,\ldots,9$).
Угорі наведено результати для $k=1$,
унизу --- для $k=5$.
Ліворуч: відносний зсув медіани
часу прибуття $b_c$;
від'ємні значення означають,
що \TS{} прогнозує раніше прибуття.
Праворуч: нормована відстань
Вассерштайна $\eW$.
Для кожної конфігурації використано
400 парних запусків. Невизначеність
не показано; результати для
термінальних прибуттів наведено
на рисунку~\ref{fig:accuracy_map}.
}
\label{fig:chain_heatmap}
\end{figure}


\subsection{Вплив розміру спільнот
і кількості міжспільнотних ребер}
\label{sec:res_size}

Попередні експерименти показали,
що внутрішня структура впливає
на розбіжності між моделями.
Однак однакова середня довжина
найкоротших шляхів не обов'язково
означає однакову точність \TS{}.

Щоб перевірити вплив інших
характеристик мережі, порівнюємо
час досягнення останньої спільноти
$T_9$ для різних розмірів $n$,
значень ASPL та кількостей
міжспільнотних ребер $k$
(рисунок~\ref{fig:accuracy_map}).

У конфігураціях із короткими
внутрішніми шляхами
($\mathrm{ASPL}\leq2$)
розбіжності переважно були невеликими.
Для 14 із 16 таких конфігурацій
перестановочний тест не виявив
статистично значущої відмінності
між розподілами часу прибуття.

За $\mathrm{ASPL}\geq3$
розбіжності ставали помітнішими.
У 31 із 32 конфігурацій
перестановочний тест виявив
відмінність між розподілами
($p<0.01$). Ці результати
отримані без поправки
на множинні порівняння.

Водночас величина розбіжності
істотно залежала від розміру
спільнот. За фіксованого
ASPL $\geq3$ нормована
відстань Вассерштайна
монотонно зменшувалася
зі збільшенням $n$
для обох значень $k$.

Наприклад, за ASPL 6 і $k=1$
вона становила 40.9\%
для $n=50$, але лише 13.6\%
для $n=1000$.
За $k=5$ відповідні значення
дорівнювали 32.6\% та 10.8\%.

Збільшення кількості
міжспільнотних ребер від одного
до п'яти також зазвичай
зменшувало розбіжності
в розріджених спільнотах,
хоча не усувало їх повністю.

\begin{figure}[tbp]
\centering
% BEGIN EDITABLE FIGURE accuracy
\begin{tikzpicture}
\begin{groupplot}[paper axis,group style={group size=2 by 2,horizontal sep=3.4cm,vertical sep=1.7cm},width=.31\linewidth,height=4.25cm,scale only axis]
\nextgroupplot[title={(a) $k=1$: median shift},xmin=-.5,xmax=3.5,ymin=-.5,ymax=5.5,xtick={0,1,2,3},xticklabels={50,100,250,1000},ytick={0,1,2,3,4,5},yticklabels={1,2,3,4,5,6},xlabel={Community size $n$},ylabel={Target internal ASPL},colormap name=paperdiv,point meta min=-40.307303343994505,point meta max=40.307303343994505,colorbar,colorbar style={width=3mm,ylabel={$b_9$ (\%)},ylabel style={font=\scriptsize},scaled y ticks=false},grid=none,tick align=outside]
\addplot[matrix plot*,mesh/cols=4,point meta=explicit] table[x=x,y=y,meta=value]{figures_pgfplots/accuracy_k1_median.dat};
\pgfplotstableread{figures_pgfplots/accuracy_k1_median.dat}\heatmapdata
\pgfplotstableforeachcolumnelement{value}\of\heatmapdata\as\cellvalue{
\pgfplotstablegetelem{\pgfplotstablerow}{x}\of\heatmapdata\edef\cellx{\pgfplotsretval}
\pgfplotstablegetelem{\pgfplotstablerow}{y}\of\heatmapdata\edef\celly{\pgfplotsretval}
\pgfplotstablegetelem{\pgfplotstablerow}{hatch}\of\heatmapdata
\ifnum\pgfplotsretval=1\relax\edef\drawcellhatch{\noexpand\path[pattern=north east lines,pattern color=gray!45] (axis cs:{\cellx-.5},{\celly-.5}) rectangle (axis cs:{\cellx+.5},{\celly+.5});}\drawcellhatch\fi
\pgfmathparse{abs(\cellvalue)>23.378235939516813}
\ifdim\pgfmathresult pt>0pt\def\celltextcolor{white}\else\def\celltextcolor{black}\fi
\edef\drawcelllabel{\noexpand\node[font=\noexpand\scriptsize,text=\celltextcolor] at (axis cs:\cellx,\celly) {\noexpand\pgfmathprintnumber[fixed,precision=1,fixed zerofill,showpos=true]{\cellvalue}};}\drawcelllabel
}
\draw[white,line width=.4pt] (axis cs:0.5,-.5) -- (axis cs:0.5,5.5);
\draw[white,line width=.4pt] (axis cs:1.5,-.5) -- (axis cs:1.5,5.5);
\draw[white,line width=.4pt] (axis cs:2.5,-.5) -- (axis cs:2.5,5.5);
\draw[white,line width=.4pt] (axis cs:-.5,0.5) -- (axis cs:3.5,0.5);
\draw[white,line width=.4pt] (axis cs:-.5,1.5) -- (axis cs:3.5,1.5);
\draw[white,line width=.4pt] (axis cs:-.5,2.5) -- (axis cs:3.5,2.5);
\draw[white,line width=.4pt] (axis cs:-.5,3.5) -- (axis cs:3.5,3.5);
\draw[white,line width=.4pt] (axis cs:-.5,4.5) -- (axis cs:3.5,4.5);
\nextgroupplot[title={(b) $k=1$: Wasserstein distance},xmin=-.5,xmax=3.5,ymin=-.5,ymax=5.5,xtick={0,1,2,3},xticklabels={50,100,250,1000},ytick={0,1,2,3,4,5},yticklabels={1,2,3,4,5,6},xlabel={Community size $n$},ylabel={Target internal ASPL},colormap name=paperseq,point meta min=0,point meta max=40.88324313976721,colorbar,colorbar style={width=3mm,ylabel={$e_W$ (\%)},ylabel style={font=\scriptsize},scaled y ticks=false},grid=none,tick align=outside]
\addplot[matrix plot*,mesh/cols=4,point meta=explicit] table[x=x,y=y,meta=value]{figures_pgfplots/accuracy_k1_w1.dat};
\pgfplotstableread{figures_pgfplots/accuracy_k1_w1.dat}\heatmapdata
\pgfplotstableforeachcolumnelement{value}\of\heatmapdata\as\cellvalue{
\pgfplotstablegetelem{\pgfplotstablerow}{x}\of\heatmapdata\edef\cellx{\pgfplotsretval}
\pgfplotstablegetelem{\pgfplotstablerow}{y}\of\heatmapdata\edef\celly{\pgfplotsretval}
\pgfplotstablegetelem{\pgfplotstablerow}{hatch}\of\heatmapdata
\ifnum\pgfplotsretval=1\relax\edef\drawcellhatch{\noexpand\path[pattern=north east lines,pattern color=gray!45] (axis cs:{\cellx-.5},{\celly-.5}) rectangle (axis cs:{\cellx+.5},{\celly+.5});}\drawcellhatch\fi
\pgfmathparse{abs(\cellvalue)>23.712281021064982}
\ifdim\pgfmathresult pt>0pt\def\celltextcolor{white}\else\def\celltextcolor{black}\fi
\edef\drawcelllabel{\noexpand\node[font=\noexpand\scriptsize,text=\celltextcolor] at (axis cs:\cellx,\celly) {\noexpand\pgfmathprintnumber[fixed,precision=1,fixed zerofill,showpos=false]{\cellvalue}};}\drawcelllabel
}
\draw[white,line width=.4pt] (axis cs:0.5,-.5) -- (axis cs:0.5,5.5);
\draw[white,line width=.4pt] (axis cs:1.5,-.5) -- (axis cs:1.5,5.5);
\draw[white,line width=.4pt] (axis cs:2.5,-.5) -- (axis cs:2.5,5.5);
\draw[white,line width=.4pt] (axis cs:-.5,0.5) -- (axis cs:3.5,0.5);
\draw[white,line width=.4pt] (axis cs:-.5,1.5) -- (axis cs:3.5,1.5);
\draw[white,line width=.4pt] (axis cs:-.5,2.5) -- (axis cs:3.5,2.5);
\draw[white,line width=.4pt] (axis cs:-.5,3.5) -- (axis cs:3.5,3.5);
\draw[white,line width=.4pt] (axis cs:-.5,4.5) -- (axis cs:3.5,4.5);
\nextgroupplot[title={(c) $k=5$: median shift},xmin=-.5,xmax=3.5,ymin=-.5,ymax=5.5,xtick={0,1,2,3},xticklabels={50,100,250,1000},ytick={0,1,2,3,4,5},yticklabels={1,2,3,4,5,6},xlabel={Community size $n$},ylabel={Target internal ASPL},colormap name=paperdiv,point meta min=-40.307303343994505,point meta max=40.307303343994505,colorbar=false,colorbar style={width=3mm,ylabel={$b_9$ (\%)},ylabel style={font=\scriptsize},scaled y ticks=false},grid=none,tick align=outside]
\addplot[matrix plot*,mesh/cols=4,point meta=explicit] table[x=x,y=y,meta=value]{figures_pgfplots/accuracy_k5_median.dat};
\pgfplotstableread{figures_pgfplots/accuracy_k5_median.dat}\heatmapdata
\pgfplotstableforeachcolumnelement{value}\of\heatmapdata\as\cellvalue{
\pgfplotstablegetelem{\pgfplotstablerow}{x}\of\heatmapdata\edef\cellx{\pgfplotsretval}
\pgfplotstablegetelem{\pgfplotstablerow}{y}\of\heatmapdata\edef\celly{\pgfplotsretval}
\pgfplotstablegetelem{\pgfplotstablerow}{hatch}\of\heatmapdata
\ifnum\pgfplotsretval=1\relax\edef\drawcellhatch{\noexpand\path[pattern=north east lines,pattern color=gray!45] (axis cs:{\cellx-.5},{\celly-.5}) rectangle (axis cs:{\cellx+.5},{\celly+.5});}\drawcellhatch\fi
\pgfmathparse{abs(\cellvalue)>23.378235939516813}
\ifdim\pgfmathresult pt>0pt\def\celltextcolor{white}\else\def\celltextcolor{black}\fi
\edef\drawcelllabel{\noexpand\node[font=\noexpand\scriptsize,text=\celltextcolor] at (axis cs:\cellx,\celly) {\noexpand\pgfmathprintnumber[fixed,precision=1,fixed zerofill,showpos=true]{\cellvalue}};}\drawcelllabel
}
\draw[white,line width=.4pt] (axis cs:0.5,-.5) -- (axis cs:0.5,5.5);
\draw[white,line width=.4pt] (axis cs:1.5,-.5) -- (axis cs:1.5,5.5);
\draw[white,line width=.4pt] (axis cs:2.5,-.5) -- (axis cs:2.5,5.5);
\draw[white,line width=.4pt] (axis cs:-.5,0.5) -- (axis cs:3.5,0.5);
\draw[white,line width=.4pt] (axis cs:-.5,1.5) -- (axis cs:3.5,1.5);
\draw[white,line width=.4pt] (axis cs:-.5,2.5) -- (axis cs:3.5,2.5);
\draw[white,line width=.4pt] (axis cs:-.5,3.5) -- (axis cs:3.5,3.5);
\draw[white,line width=.4pt] (axis cs:-.5,4.5) -- (axis cs:3.5,4.5);
\nextgroupplot[title={(d) $k=5$: Wasserstein distance},xmin=-.5,xmax=3.5,ymin=-.5,ymax=5.5,xtick={0,1,2,3},xticklabels={50,100,250,1000},ytick={0,1,2,3,4,5},yticklabels={1,2,3,4,5,6},xlabel={Community size $n$},ylabel={Target internal ASPL},colormap name=paperseq,point meta min=0,point meta max=40.88324313976721,colorbar=false,colorbar style={width=3mm,ylabel={$e_W$ (\%)},ylabel style={font=\scriptsize},scaled y ticks=false},grid=none,tick align=outside]
\addplot[matrix plot*,mesh/cols=4,point meta=explicit] table[x=x,y=y,meta=value]{figures_pgfplots/accuracy_k5_w1.dat};
\pgfplotstableread{figures_pgfplots/accuracy_k5_w1.dat}\heatmapdata
\pgfplotstableforeachcolumnelement{value}\of\heatmapdata\as\cellvalue{
\pgfplotstablegetelem{\pgfplotstablerow}{x}\of\heatmapdata\edef\cellx{\pgfplotsretval}
\pgfplotstablegetelem{\pgfplotstablerow}{y}\of\heatmapdata\edef\celly{\pgfplotsretval}
\pgfplotstablegetelem{\pgfplotstablerow}{hatch}\of\heatmapdata
\ifnum\pgfplotsretval=1\relax\edef\drawcellhatch{\noexpand\path[pattern=north east lines,pattern color=gray!45] (axis cs:{\cellx-.5},{\celly-.5}) rectangle (axis cs:{\cellx+.5},{\celly+.5});}\drawcellhatch\fi
\pgfmathparse{abs(\cellvalue)>23.712281021064982}
\ifdim\pgfmathresult pt>0pt\def\celltextcolor{white}\else\def\celltextcolor{black}\fi
\edef\drawcelllabel{\noexpand\node[font=\noexpand\scriptsize,text=\celltextcolor] at (axis cs:\cellx,\celly) {\noexpand\pgfmathprintnumber[fixed,precision=1,fixed zerofill,showpos=false]{\cellvalue}};}\drawcelllabel
}
\draw[white,line width=.4pt] (axis cs:0.5,-.5) -- (axis cs:0.5,5.5);
\draw[white,line width=.4pt] (axis cs:1.5,-.5) -- (axis cs:1.5,5.5);
\draw[white,line width=.4pt] (axis cs:2.5,-.5) -- (axis cs:2.5,5.5);
\draw[white,line width=.4pt] (axis cs:-.5,0.5) -- (axis cs:3.5,0.5);
\draw[white,line width=.4pt] (axis cs:-.5,1.5) -- (axis cs:3.5,1.5);
\draw[white,line width=.4pt] (axis cs:-.5,2.5) -- (axis cs:3.5,2.5);
\draw[white,line width=.4pt] (axis cs:-.5,3.5) -- (axis cs:3.5,3.5);
\draw[white,line width=.4pt] (axis cs:-.5,4.5) -- (axis cs:3.5,4.5);
\end{groupplot}
\end{tikzpicture}
% END EDITABLE FIGURE accuracy
\caption{%
Розбіжності термінального часу
прибуття $T_9$ для 48 конфігурацій
ланцюгів із десяти спільнот.
Ліві панелі показують відносний зсув медіани
$b_9$ (\%), праві --- нормовану відстань
Вассерштайна $\eW$ (\%). Верхній ряд панелей
відповідає $k=1$, нижній --- $k=5$.
У кожній панелі рядки комірок відповідають
цільовій ASPL, а стовпці --- розміру спільноти $n$.
Колір і число в комірці відображають один і той самий
показник; кольорові шкали спільні для двох значень $k$
окремо для кожного показника.
Від'ємний медіанний зсув означає раніше прибуття у \TS{}.
Штрихування лише на правих панелях позначає
конфігурації, для яких перестановочний тест відстані
Вассерштайна не виявив відмінності ($p\geq0.05$).
Перевірки проведено без поправки
на множинні порівняння.
Для кожної конфігурації
використано 400 парних запусків.
}
\label{fig:accuracy_map}
\end{figure}

Частка міжспільнотних переходів
\TS{}, які відбувалися до зараження
першого справжнього моста в джерелі,
також істотно залежала від структури:
від 0.1--2.3\% за ASPL 1
до 31--51\% за ASPL 6.

Отже, ASPL є корисною характеристикою
для оцінювання точності TwoScale:
зі збільшенням довжини внутрішніх
шляхів розбіжності між моделями
зазвичай зростають. Проте за однакового
ASPL їхня величина може суттєво
відрізнятися залежно від розміру
спільноти та кількості міжспільнотних
ребер. Це свідчить про те, що ASPL
сама по собі недостатня для оцінювання
точності, і мотивує пошук характеристики,
яка також враховує часову динаміку
внутрішнього поширення.


\subsection{Зв'язок розбіжностей
зі співвідношенням часових масштабів}
\label{sec:res_descriptor}

Отримані результати свідчать,
що розбіжності між моделями
можуть залежати від того,
наскільки швидко інфекція
поширюється всередині спільноти
порівняно з міжспільнотним
передаванням.

Щоб перевірити цей зв'язок,
використовуємо показник
$\rho=\beta_{\mathrm{out}}\sigma_T$,
визначений у розділі 5.3.
Він характеризує розкид
внутрішніх часів зараження
відносно середнього часу
передавання через один
міжспільнотний канал.

На рисунку~\ref{fig:descriptor}(a,b)
показано залежності нормованої
відстані Вассерштайна $\eW$
та зсуву медіани $b_9$
від $\rho$ для всіх
48 конфігурацій ланцюгів.

Показник $\rho$ мав сильніший
ранговий зв'язок із $\eW$,
ніж цільова ASPL.
Кореляція Спірмена між
$\rho$ та $\eW$ становила
0.91 для всіх конфігурацій
і 0.98 для конфігурацій
із $\mathrm{ASPL}\geq3$.
Відповідні значення
для ASPL становили 0.82 і 0.66.

За $\rho<0.25$ відстань $\eW$
не перевищувала 4.5\%
у всіх 18 конфігураціях.
Для 15 із них перестановочний
тест не виявив відмінності
між розподілами.
Натомість за $\rho\geq0.25$
спостережувана відстань
перевищувала перестановочний
поріг у всіх 30 конфігураціях.

Для $\rho\geq0.3$
одержані результати наближено
описувалися емпіричним
співвідношенням
\begin{equation}
\eW \approx 15\rho\,\%.
\end{equation}
У цій області також
спостерігалося
$b_9\approx-\eW$:
відстань між розподілами
переважно відповідала
зміщенню часу прибуття
в бік раніших значень
у \TS{}.

Таким чином, результати
підтримують пояснення
через співвідношення
часових масштабів.
Коли розкид внутрішніх
часів зараження малий
порівняно з очікуванням
міжспільнотного передавання,
агрегація вузлів-мостів
має порівняно невеликий
вплив на часові розподіли.

Водночас показник $\rho$
не є універсальною формулою
похибки. Наведені порогові
значення та емпіричне
наближення отримані
за тими самими конфігураціями,
на яких оцінювали розбіжності.
Їхню здатність прогнозувати
результати на незалежних
мережах тут не перевірено.

\begin{figure}[tbp]
\centering
% BEGIN EDITABLE FIGURE descriptor
\begin{tikzpicture}
\begin{groupplot}[paper axis,group style={group size=3 by 1,horizontal sep=1.7cm},width=.24\linewidth,height=4.1cm,scale only axis]
\nextgroupplot[xmode=log,xlabel={$\rho=\beta\sigma_T$},ylabel={$e_W$ of $T_9$ (\%)},title={(a) Distributional discrepancy}]
\addplot[Size50,mark=*,mark size=1.7pt,mark options={fill=Size50}] table[x=plotrho,y=rel_w1_pct]{figures_pgfplots/descriptor_n50_k1.dat};
\addplot[gray!65,only marks,mark=-,mark size=2pt] table[x=plotrho,y=rel_w1_null95_pct]{figures_pgfplots/descriptor_n50_k1.dat};
\addplot[Size50,mark=triangle,mark size=1.7pt,mark options={fill=white}] table[x=plotrho,y=rel_w1_pct]{figures_pgfplots/descriptor_n50_k5.dat};
\addplot[gray!65,only marks,mark=-,mark size=2pt] table[x=plotrho,y=rel_w1_null95_pct]{figures_pgfplots/descriptor_n50_k5.dat};
\addplot[Size100,mark=*,mark size=1.7pt,mark options={fill=Size100}] table[x=plotrho,y=rel_w1_pct]{figures_pgfplots/descriptor_n100_k1.dat};
\addplot[gray!65,only marks,mark=-,mark size=2pt] table[x=plotrho,y=rel_w1_null95_pct]{figures_pgfplots/descriptor_n100_k1.dat};
\addplot[Size100,mark=triangle,mark size=1.7pt,mark options={fill=white}] table[x=plotrho,y=rel_w1_pct]{figures_pgfplots/descriptor_n100_k5.dat};
\addplot[gray!65,only marks,mark=-,mark size=2pt] table[x=plotrho,y=rel_w1_null95_pct]{figures_pgfplots/descriptor_n100_k5.dat};
\addplot[Size250,mark=*,mark size=1.7pt,mark options={fill=Size250}] table[x=plotrho,y=rel_w1_pct]{figures_pgfplots/descriptor_n250_k1.dat};
\addplot[gray!65,only marks,mark=-,mark size=2pt] table[x=plotrho,y=rel_w1_null95_pct]{figures_pgfplots/descriptor_n250_k1.dat};
\addplot[Size250,mark=triangle,mark size=1.7pt,mark options={fill=white}] table[x=plotrho,y=rel_w1_pct]{figures_pgfplots/descriptor_n250_k5.dat};
\addplot[gray!65,only marks,mark=-,mark size=2pt] table[x=plotrho,y=rel_w1_null95_pct]{figures_pgfplots/descriptor_n250_k5.dat};
\addplot[Size1000,mark=*,mark size=1.7pt,mark options={fill=Size1000}] table[x=plotrho,y=rel_w1_pct]{figures_pgfplots/descriptor_n1000_k1.dat};
\addplot[gray!65,only marks,mark=-,mark size=2pt] table[x=plotrho,y=rel_w1_null95_pct]{figures_pgfplots/descriptor_n1000_k1.dat};
\addplot[Size1000,mark=triangle,mark size=1.7pt,mark options={fill=white}] table[x=plotrho,y=rel_w1_pct]{figures_pgfplots/descriptor_n1000_k5.dat};
\addplot[gray!65,only marks,mark=-,mark size=2pt] table[x=plotrho,y=rel_w1_null95_pct]{figures_pgfplots/descriptor_n1000_k5.dat};
\nextgroupplot[xmode=log,xlabel={$\rho=\beta\sigma_T$},ylabel={Median shift $b_9$ (\%)},title={(b) Signed median shift}]
\addplot[gray,no marks] coordinates {(0.0018131616755314697,0) (2.7344524800897014,0)};
\addplot[Size50,mark=*,mark size=1.7pt,mark options={fill=Size50}] table[x=plotrho,y=median_shift_pct]{figures_pgfplots/descriptor_n50_k1.dat};
\addplot[Size50,mark=triangle,mark size=1.7pt,mark options={fill=white}] table[x=plotrho,y=median_shift_pct]{figures_pgfplots/descriptor_n50_k5.dat};
\addplot[Size100,mark=*,mark size=1.7pt,mark options={fill=Size100}] table[x=plotrho,y=median_shift_pct]{figures_pgfplots/descriptor_n100_k1.dat};
\addplot[Size100,mark=triangle,mark size=1.7pt,mark options={fill=white}] table[x=plotrho,y=median_shift_pct]{figures_pgfplots/descriptor_n100_k5.dat};
\addplot[Size250,mark=*,mark size=1.7pt,mark options={fill=Size250}] table[x=plotrho,y=median_shift_pct]{figures_pgfplots/descriptor_n250_k1.dat};
\addplot[Size250,mark=triangle,mark size=1.7pt,mark options={fill=white}] table[x=plotrho,y=median_shift_pct]{figures_pgfplots/descriptor_n250_k5.dat};
\addplot[Size1000,mark=*,mark size=1.7pt,mark options={fill=Size1000}] table[x=plotrho,y=median_shift_pct]{figures_pgfplots/descriptor_n1000_k1.dat};
\addplot[Size1000,mark=triangle,mark size=1.7pt,mark options={fill=white}] table[x=plotrho,y=median_shift_pct]{figures_pgfplots/descriptor_n1000_k5.dat};
\nextgroupplot[title={(c) First vs later crossings},xmin=-47,xmax=5,ymin=-47,ymax=5,xtick={-40,-20,0},ytick={-40,-20,0},xlabel={First crossing: mean shift (\%)},ylabel={Crossings 2--9: mean shift (\%)}]
\addplot[gray,no marks] coordinates {(-47,-47) (5,5)};
\addplot[only marks,Size50,mark=*,mark size=1.7pt,mark options={fill=Size50},error bars/.cd,x dir=both,x explicit,error bar style={line width=.3pt},error mark options={mark size=0pt}] table[x=first_cross_shift_pct,y=later_cross_shift_pct,x error minus=errlo,x error plus=errhi]{figures_pgfplots/descriptor_cross_n50_k1.dat};
\addplot[only marks,Size50,mark=triangle,mark size=1.7pt,mark options={fill=white},error bars/.cd,x dir=both,x explicit,error bar style={line width=.3pt},error mark options={mark size=0pt}] table[x=first_cross_shift_pct,y=later_cross_shift_pct,x error minus=errlo,x error plus=errhi]{figures_pgfplots/descriptor_cross_n50_k5.dat};
\addplot[only marks,Size100,mark=*,mark size=1.7pt,mark options={fill=Size100},error bars/.cd,x dir=both,x explicit,error bar style={line width=.3pt},error mark options={mark size=0pt}] table[x=first_cross_shift_pct,y=later_cross_shift_pct,x error minus=errlo,x error plus=errhi]{figures_pgfplots/descriptor_cross_n100_k1.dat};
\addplot[only marks,Size100,mark=triangle,mark size=1.7pt,mark options={fill=white},error bars/.cd,x dir=both,x explicit,error bar style={line width=.3pt},error mark options={mark size=0pt}] table[x=first_cross_shift_pct,y=later_cross_shift_pct,x error minus=errlo,x error plus=errhi]{figures_pgfplots/descriptor_cross_n100_k5.dat};
\addplot[only marks,Size250,mark=*,mark size=1.7pt,mark options={fill=Size250},error bars/.cd,x dir=both,x explicit,error bar style={line width=.3pt},error mark options={mark size=0pt}] table[x=first_cross_shift_pct,y=later_cross_shift_pct,x error minus=errlo,x error plus=errhi]{figures_pgfplots/descriptor_cross_n250_k1.dat};
\addplot[only marks,Size250,mark=triangle,mark size=1.7pt,mark options={fill=white},error bars/.cd,x dir=both,x explicit,error bar style={line width=.3pt},error mark options={mark size=0pt}] table[x=first_cross_shift_pct,y=later_cross_shift_pct,x error minus=errlo,x error plus=errhi]{figures_pgfplots/descriptor_cross_n250_k5.dat};
\addplot[only marks,Size1000,mark=*,mark size=1.7pt,mark options={fill=Size1000},error bars/.cd,x dir=both,x explicit,error bar style={line width=.3pt},error mark options={mark size=0pt}] table[x=first_cross_shift_pct,y=later_cross_shift_pct,x error minus=errlo,x error plus=errhi]{figures_pgfplots/descriptor_cross_n1000_k1.dat};
\addplot[only marks,Size1000,mark=triangle,mark size=1.7pt,mark options={fill=white},error bars/.cd,x dir=both,x explicit,error bar style={line width=.3pt},error mark options={mark size=0pt}] table[x=first_cross_shift_pct,y=later_cross_shift_pct,x error minus=errlo,x error plus=errhi]{figures_pgfplots/descriptor_cross_n1000_k5.dat};
\node[anchor=north west,font=\tiny,align=left,text=gray] at (axis cs:-45,3) {ASPL $\geq3$\\Below diagonal: later\\crossings more biased};
\end{groupplot}
\node[anchor=north,font=\scriptsize] at ([yshift=-1.25cm]group c2r1.south) {\textcolor{Size50}{$\blacksquare$} $n=50$\quad\textcolor{Size100}{$\blacksquare$} $n=100$\quad\textcolor{Size250}{$\blacksquare$} $n=250$\quad\textcolor{Size1000}{$\blacksquare$} $n=1000$\quad$\bullet$ $k=1$\quad$\triangle$ $k=5$\quad\textcolor{gray}{---} null 95\%};
\end{tikzpicture}
% END EDITABLE FIGURE descriptor
\caption{%
Зв'язок розбіжностей між \Full{}
і \TS{} зі співвідношенням
часових масштабів
$\rho=\beta_{\mathrm{out}}\sigma_T$
для 48 конфігурацій ланцюгів.
(a) Нормована відстань Вассерштайна
$\eW$ для термінального часу
прибуття $T_9$.
(b) Відносний зсув його медіани $b_9$.
Кольори позначають розмір спільноти;
зафарбовані кола відповідають $k=1$,
порожні трикутники --- $k=5$.
Лінії сполучають конфігурації
з ASPL від 1 до 6 за фіксованих
$n$ і $k$. Сірі позначки в (a)
показують 95-й процентиль
перестановочного розподілу
відстані для відповідної конфігурації.
(c) Порівняння відносних зсувів
середньої тривалості першого
переходу та переходів 2--9
для $\mathrm{ASPL}\geq3$.
Відрізки похибок --- 95\%
блокові бутстреп-інтервали
зсуву першого переходу.
}
\label{fig:descriptor}
\end{figure}


\subsection{Відмінності між першим
і наступними міжспільнотними переходами}

Попередні результати показали,
що відносна розбіжність часу
прибуття часто стабілізується
вздовж ланцюга.
Проте це не означає,
що перший і наступні переходи
мають однакове зміщення.

Перший перехід відбувається
з початкової спільноти,
яка не отримувала імпортованих
заражень. Тому на його час
впливають насамперед
представлення інтерфейсу
джерела та чисельне
визначення моменту експорту.

На наступних переходах
можуть проявлятися
додаткові наслідки
агрегованого моделювання:
випадкове розміщення
імпортованих заражень
і можливість повторних
імпортів через інтерфейси,
фізичні канали яких
уже не можуть спричинити
нові зараження.

Щоб перевірити, чи відрізняється
зміщення цих етапів,
порівнюємо відносні зсуви
середніх тривалостей
першого переходу $T_1$
та наступних переходів
$T_9-T_1$
(рисунок~\ref{fig:descriptor}(c)).

Для $k=1$ у 15 із 16
конфігурацій із
$\mathrm{ASPL}\geq3$
наступні переходи мали
сильніше від'ємне зміщення,
ніж перший.
У дев'яти з них
блокові довірчі інтервали
для різниці зсувів
не містили нуль.

Наприклад, за $n=50$
і ASPL 6 зсув середньої
тривалості першого переходу
становив $-27.5\%$,
тоді як для наступних
переходів --- $-42.5\%$.

За $k=5$ ця відмінність
була менш вираженою:
сильніше від'ємне
зміщення наступних
переходів спостерігалося
у 14 із 16 конфігурацій,
причому для чотирьох
із них довірчий інтервал
різниці не містив нуль.
Перевірки виконано
без поправки на множинні
порівняння.

Ця залежність сумісна
з припущенням, що вузькі
інтерфейси чутливіші
до надлишкових імпортів.
За одного фізичного ребра
можливе лише одне
зараження його цільового
вузла, тоді як агрегована
інтенсивність може
спричиняти подальші імпорти.
За більшої кількості
міжспільнотних ребер
це обмеження є менш
жорстким.

Водночас сильніше зміщення
наступних переходів
не суперечить повільному
зростанню або стабілізації
відносної розбіжності
кумулятивного часу прибуття.

Для середніх часів маємо
\begin{equation}
\mathbb{E}[T_9]
=
\mathbb{E}[T_1]
+
\mathbb{E}[T_9-T_1].
\end{equation}
Тому відносний зсув
середнього $T_9$
є зваженим середнім
відносних зсувів
цих двох компонентів,
де ваги визначаються
їхніми середніми
тривалостями у \Full{}.

Отже, абсолютне випередження
може накопичуватися разом
зі збільшенням часу поширення,
тоді як відносна розбіжність
наближається до характерного
зміщення наступних переходів.
Це пояснення безпосередньо
стосується середніх часів;
поведінку медіан на рисунку
\ref{fig:chain_heatmap} воно
не встановлює математично.

Загалом результати вказують
на відмінності між першим
і наступними міжспільнотними
переходами. Однак саме
порівняння їхніх тривалостей
не дозволяє окремо визначити
внесок випадкового розміщення
імпортів, надлишкових
заражень і чисельної
синхронізації.
Для цього потрібні додаткові
контрольні експерименти,
які ізолюють кожен
із зазначених механізмів.







